% §1.6 — La communication sans fil et la radiodiffusion (~1895–1930)
% VERSION : v1 — 27 avril 2026
% STRUCTURE : quatre temps + bilan
%   (1) Ouverture : la libération du fil (Hertz → Marconi, cluster VdK, brevet 1897)
%   (2) Le monopole Marconi et le modèle maritime (3 piliers, Poldhu, Telefunken, WWI → RCA)
%   (3) La puissance dans le signal (Aitken : spark < arc, Navy $1.5M, câbles vulnérables)
%   (4) Le spectre et le broadcasting (Hazlett 1990, Adena et al. 2015, Gentzkow 2006)
%   Bilan : tableau D1–D10 + 0-TER, trois hypothèses abductives
% SOURCES PRINCIPALES : VdK wireless_engine, Aitken 1985 (ch. 2–3, 5),
%   Hazlett 1990 (JLE), Adena et al. 2015 (QJE/WZB), Gentzkow 2006 (QJE),
%   Jolly 1972, Raboy 2016, Douglas 1987
% =============================================================================

\section{La communication sans fil et la radiodiffusion (~1895--1930)}\label{sec:wireless}

Le télégraphe et le téléphone transmettaient de l'information à travers un fil de cuivre~; le wireless transmet à travers l'air. Le changement paraît simple, mais il modifie la structure même du couplage entre information et matière. Dans les cas précédents, la matérialité du réseau résidait dans le substrat~: le cuivre, les poteaux, la gutta-percha, les isolateurs. Le coût de transmission croissait avec la distance parce que le fil croissait avec la distance. Avec le wireless, la matérialité se déplace du substrat au signal~: ce n'est plus la longueur du conducteur qui détermine la portée, c'est la puissance de l'émetteur. L'infrastructure n'est plus un réseau qu'on pose~; c'est une énergie qu'on émet. Ce déplacement, qui paraît libérateur parce qu'il affranchit la transmission du câble, introduit en réalité un rapport qualitativement nouveau à l'énergie~: pour la première fois dans le parcours de ce chapitre, le coût d'exploitation énergétique de la transmission est continu et croissant avec l'ambition du service.

\dimmark{D10}{+}%
La démonstration par Heinrich Hertz, entre 1886 et 1888, de l'existence des ondes électromagnétiques prédites par Maxwell ne déboucha pas sur une application~: Hertz lui-même ne voyait aucun usage pratique à sa découverte \parencite{jolly_marconi_1972}. Ce qui suivit, entre 1890 et 1896, fut un cluster d'innovation au sens de \textcite{kooij_information_2022}~: Lodge en Angleterre, Branly en France, Righi en Italie, Popov en Russie, Tesla aux États-Unis, Bose en Inde travaillèrent indépendamment sur la détection des ondes hertziennes, sans qu'aucun d'entre eux ne construise un système de communication opérationnel \parencite{jolly_marconi_1972, raboy_marconi_2016}. Le passage du phénomène de laboratoire au système technique fut l'œuvre de Guglielmo Marconi, un jeune Italien de vingt et un ans sans formation universitaire, qui assembla les composants disponibles (oscillateur de Righi, cohéreur de Branly, antenne, terre) en un dispositif capable de transmettre un signal sur plusieurs kilomètres. Marconi n'inventa aucun composant~; il inventa le système, de la même manière qu'Edison n'avait pas inventé la lampe à incandescence mais le réseau électrique (§\,\ref{sec:electrification}). Le brevet britannique n\textsuperscript{o}~12,039, déposé en juin~1897, couvrait non pas un appareil mais un procédé de transmission de signaux électriques sans fil conducteur \parencite{jolly_marconi_1972}.


\subsection{Le monopole Marconi et le modèle maritime}\label{subsec:marconi}

\dimmark{D2/D7}{*}%
La stratégie commerciale de Marconi reposait sur trois piliers dont la combinaison produisit, pendant deux décennies, un monopole de fait sur la communication sans fil maritime. Le premier pilier était le portefeuille de brevets~: la \emph{Wireless Telegraph and Signal Company}, fondée en juillet~1897 avec un capital de 100\,000~livres sterling, investit environ 75\,000~livres dans l'acquisition et le dépôt de brevets avant~1900 \parencite{kooij_information_2022}. Le second était l'internationalisation~: plutôt que de vendre des licences à des opérateurs étrangers, Marconi créa des filiales nationales sous son contrôle (la \emph{Marconi International Marine Communication Company}, MIMCC, en~1900~; des filiales en Belgique, au Canada, en Argentine) et imposa à chacune l'obligation d'utiliser exclusivement le matériel Marconi \parencite{jolly_marconi_1972, kooij_information_2022}. Le troisième pilier était le leasing~: les équipements n'étaient pas vendus aux armateurs mais loués, accompagnés d'un opérateur Marconi embarqué, ce qui liait le client au fournisseur par un contrat de service continu\footnote{Le parallèle avec le modèle IBM (§\,\ref{sec:mechanography}) est frappant~: location de l'équipement, vente exclusive du consommable (les cartes chez IBM, le service d'opérateur chez Marconi), portefeuille de brevets comme barrière à l'entrée. Les deux monopoles procèdent du même triptyque, appliqué à des technologies sans rapport.}. Le refus d'interopérabilité~: Marconi interdit à ses opérateurs de communiquer avec les stations équipées de matériel concurrent, ce qui provoqua un conflit diplomatique à la conférence radiotélégraphique de Berlin en~1903 et conduisit à la Convention radiotélégraphique internationale de~1906, première tentative de régulation du spectre \parencite{jolly_marconi_1972, headrick_invisible_1991}.

\dimmark{0-TER}{+}%
La traversée transatlantique du 12~décembre~1901, lorsque Marconi reçut à St.~John's (Terre-Neuve) un signal émis depuis Poldhu (Cornouailles), marqua un seuil dans le rapport entre transmission et énergie. L'émetteur de Poldhu, un spark transmitter alimenté par un alternateur de 25~kilowatts, consommait une puissance comparable à celle d'une petite usine \parencite{jolly_marconi_1972, aitken_continuous_1985}. À seuil de réception donné, la portée variait comme la racine de la puissance d'émission~: doubler la portée exigeait de quadrupler la puissance, ce qui inscrivait la communication sans fil dans une logique de coût croissant radicalement différente de celle du télégraphe filaire.

\dimmark{D2/D8}{+}%
La concurrence vint de l'État. En Allemagne, l'administration impériale soutint la création de Telefunken (1903), société commune de Siemens et AEG, pour contrer le monopole Marconi. Telefunken bénéficia de contrats militaires et d'une politique de brevets agressive \parencite{kooij_information_2022, headrick_invisible_1991}. Le modèle institutionnel divergeait~: là où Marconi opérait comme un monopole privé international, Telefunken fonctionnait comme un instrument de politique industrielle. Les ventes cumulées de Marconi n'atteignaient que 6\,092~livres sterling en~1900 \parencite{kooij_information_2022}~; c'est la promesse stratégique, pas le marché, qui finançait l'expansion. La Première Guerre mondiale trancha le conflit~: les marines britannique et américaine prirent le contrôle des stations radio, et le monopole Marconi fut liquidé aux États-Unis en~1919 par la création de la Radio Corporation of America (RCA), orchestrée par la Navy comme \og \emph{chosen instrument}\fg{} de la politique de communication américaine \parencite{aitken_continuous_1985, kooij_information_2022}.

\dimmark{D5/D2}{*}%
\emergence{D2~--- Le monopole Marconi procède d'un quatrième mécanisme, distinct des rendements de réseau (télégraphe), du verrouillage du consommable (tabulatrice) et de l'intégration verticale (téléphone)~: le triptyque brevets + leasing + refus d'interopérabilité. Mais c'est l'État qui le défait (RCA)~: le wireless est la première technologie de communication où l'expropriation étatique en temps de guerre, et non la concurrence marchande, met fin au monopole privé. \questionch{2}{D2~--- L'État comme destructeur de monopole et créateur simultané d'un autre (RCA)~: ce mécanisme se retrouve-t-il dans les cas ultérieurs~?}}


\subsection{La puissance dans le signal}\label{subsec:puissance}

\dimmark{D8/0-TER}{*}%
L'exigence de puissance croissante posa un problème technique que les archives exploitées par \textcite{aitken_continuous_1985} permettent de documenter avec précision. En~1909, la Navy américaine publia un appel d'offres pour sa station d'Arlington (Virginie), premier élément d'un réseau de stations puissantes devant couvrir le Pacifique (Canal de Panama, Californie, Hawaï, Guam, Samoa, Philippines). Le Congrès alloua 1~million de dollars en~1911, puis augmenta l'enveloppe à 1,5~million. Le cahier des charges exigeait une portée de 3\,000~miles en toutes conditions.

La \emph{National Electric Signaling Company} (NESCO) de Fessenden remporta le contrat avec un spark transmitter synchrone rotatif de 100~kilowatts. La machine échoua à atteindre les spécifications. Ses quarante-huit électrodes en cuivre tournant à 1\,250~tours par minute dissipaient une part considérable de leur énergie en harmoniques inutiles \parencite{aitken_continuous_1985}. Un troisième émetteur, installé discrètement à Arlington par Cyril Elwell de la Federal Telegraph Company, était un arc de 30~kilowatts seulement, soit moins d'un tiers de la puissance du spark Fessenden. Il fit mieux.

\dimmark{0-TER}{*}%
Le fait que 100~kilowatts de spark produisent une portée inférieure à 30~kilowatts d'arc est le résultat 0-TER central de ce cas. L'explication est physique~: le spark dissipe son énergie en un spectre large d'harmoniques, tandis que l'arc concentre la sienne sur une seule fréquence, ce qui maximise le transfert d'énergie vers le récepteur\footnote{L'amiral Cone, chef du Bureau of Steam Engineering, objecta à Elwell qu'un ampère était un ampère et qu'on ne pouvait pas le changer (\og \emph{an ampere is an ampere and you cannot change it}\fg{}). Elwell rétorqua qu'un ampère d'onde continue valait deux ampères d'onde amortie. Cone avait tort~: c'est l'efficacité spectrale, pas la puissance brute, qui détermine la portée. L'épisode est rapporté par \textcite{aitken_continuous_1985}, ch.~3.}. La transition du spark à l'onde continue (arc, puis alternateur, puis tube à vide) est d'abord une innovation d'efficacité énergétique~: elle ne réduit pas la consommation totale mais augmente la portée par watt émis. La trajectoire de puissance qui suivit le confirme~: alternateur Alexanderson 2~kilowatts (1909), arc Federal 30~kilowatts (Arlington, 1913), arc 200~kilowatts et 500~kilowatts dans les années qui suivirent \parencite{aitken_continuous_1985}.

\dimmark{D5/0-GOV}{*}%
Le réseau de stations du Pacifique s'inscrivait dans une logique de souveraineté dont la Première Guerre mondiale allait révéler l'urgence. Le wireless répondait à une vulnérabilité stratégique que les câbles sous-marins avaient créée. \textcite{aitken_continuous_1985} documente cette vulnérabilité en détail~: l'Allemagne, coupée des câbles par la Grande-Bretagne dès~1914, ne disposait que de la station de Nauen pour communiquer avec ses diplomates et ses agents. Le télégramme Zimmerman de~1917, dans lequel l'Allemagne proposait une alliance au Mexique en échange de l'Arizona, du Nouveau-Mexique et du Texas, fut intercepté par le \emph{Naval Intelligence} britannique parce qu'il transitait par les câbles diplomatiques américains que Wilson avait mis à la disposition de l'Allemagne\footnote{L'affaire Zimmerman est rapportée par \textcite{aitken_continuous_1985}, ch.~5. Le président de Western Union, Newcomb Carlton, admit devant le Sénat en~1921 que tous les câbles entrant et sortant de Grande-Bretagne étaient remis quotidiennement au \emph{Naval Intelligence} britannique. Carlton fut assuré qu'ils n'étaient \og ni censurés ni même déchiffrés\fg{}~; les sacs de câbles étaient \og stockés dans un entrepôt de l'Amirauté et restitués le lendemain matin, non ouverts\fg{}.}. Au-delà de la diplomatie, le contrôle britannique des câbles s'étendait au commerce~: des négociants américains se plaignirent que leurs messages étaient retardés ou transmis à des concurrents britanniques par l'intermédiaire du \emph{Board of Trade} \parencite{aitken_continuous_1985}. Le wireless naissait comme technologie de souveraineté, pas de marché~: c'est la vulnérabilité des câbles qui finançait le développement de la radio, comme c'est la sécurité ferroviaire qui avait financé le déploiement du télégraphe en Angleterre (§\,\ref{sec:telegraph}). Le même pattern se reproduisait~: la demande étatique créait le marché avant que le marché n'existe.

\emergence{D5/0-GOV~--- La Navy finance le réseau Pacifique (1,5~million de dollars), exproprie le wireless pendant la guerre, crée RCA comme instrument national. La demande militaire crée le marché avant la demande commerciale~: même pattern que le Census pour Hollerith (§\,\ref{sec:mechanography}) et la Convention pour le télégraphe Chappe (§\,\ref{sec:telegraph}). \questionch{2}{D5~--- L'État comme premier client des technologies de l'information est-il un fait contingent (guerres, crises) ou structurel (les rendements croissants de réseau ne peuvent être financés que par un acteur qui ne calcule pas en termes de rentabilité immédiate)~?}}


\subsection{Le spectre et le broadcasting}\label{subsec:broadcasting}

Le troisième temps de l'histoire du wireless est celui où la technique change de nature~: de la communication point-à-point (comme le télégraphe) à la diffusion point-à-tous (le broadcasting). Deux faits nouveaux par rapport aux cas précédents émergent de cette transition.

\dimmark{D2}{*}%
Le premier fait est l'apparition d'un type de ressource sans précédent dans l'histoire des techniques de l'information. Le spectre électromagnétique est une ressource immatérielle et congestionnable~: deux émetteurs sur la même fréquence s'annulent mutuellement, comme deux trains sur la même voie se percutent. Mais, à la différence du cuivre ou du fil, le spectre n'est ni stockable, ni transportable, ni appropriable par la possession physique. C'est le premier bien commun de l'histoire de l'ICT~: ni bien privé rival (comme le cuivre du télégraphe), ni service (comme la commutation téléphonique), ni bien public pur (puisqu'il est congestionnable).

La manière dont ce bien commun fut régulé constitue l'un des épisodes les mieux documentés de l'économie de la régulation. \textcite{coase_federal_1959}, dans un article fondateur, avait soutenu que la régulation du spectre par la \emph{Federal Communications Commission} (FCC) était une erreur~: le problème de l'interférence aurait pu être résolu par la définition de droits de propriété privés sur les fréquences, comme on définit des droits de propriété sur la terre, et le marché aurait ensuite alloué ces droits à leurs usages les plus valorisés. \textcite{hazlett_rationality_1990}, dans une révision minutieuse de l'histoire législative, a montré que cette interprétation par l'erreur était elle-même erronée. Des droits de propriété privés sur le spectre existaient \emph{de facto} avant le Radio Act de~1927~: le Département du Commerce reconnaissait les transferts de licences, les stations de radio étaient achetées et vendues à des prix reflétant la valeur de la fréquence\footnote{La vente la plus spectaculaire eut lieu en septembre~1926~: la station WEAF de New York fut vendue par AT\&T à RCA pour 1~million de dollars, dont 800\,000~dollars pour le droit de fréquence seul \parencite{hazlett_rationality_1990}.}, et les stations partageant une même fréquence négociaient entre elles la répartition du temps d'antenne par des accords privés. En octobre~1926, le tribunal de l'Illinois statua, dans l'affaire \emph{Tribune Co.~v.~Oak Leaves Broadcasting Station}, que la priorité d'usage créait un droit de propriété sur la fréquence, en s'appuyant sur le droit commun des droits riverains et de la concurrence déloyale \parencite{hazlett_rationality_1990}.

\dimmark{D2/0-GOV}{*}%
Le Congrès réagit immédiatement~: en décembre~1926, les deux chambres votèrent une résolution déclarant qu'aucun droit privé sur le spectre ne serait reconnu, et exigèrent des broadcasters qu'ils signent des renonciations. Le Federal Radio Act de février~1927 créa la Federal Radio Commission et lui confia le pouvoir d'attribuer des licences selon le critère de l'\og intérêt public, de la commodité ou de la nécessité\fg{}. Ce que Hazlett montre, c'est que cette solution ne fut pas le produit d'une confusion entre la définition des droits et leur allocation~: elle fut un choix rationnel de distribution de rentes, dans lequel les grands broadcasters obtenaient la protection de leurs positions acquises (le gel de la bande à 550--1\,500~kHz, le refus d'élargir le spectre malgré la disponibilité technique), les régulateurs obtenaient un pouvoir discrétionnaire sur la vie et la mort des stations, et le Congrès obtenait un levier politique sur le contenu des émissions. La commission gela la bande AM à sa taille de~1924, utilisant moins de cinq pour cent de la capacité alors techniquement disponible \parencite{hazlett_rationality_1990}. Le \og chaos\fg{} de juillet~1926 à février~1927, invoqué par la Cour suprême comme justification de la régulation, avait été stratégiquement provoqué par le secrétaire au Commerce Herbert Hoover, qui refusa de faire appel du verdict \emph{Zenith} et de convoquer la conférence radio annuelle\footnote{Hoover avait déclaré qu'il \og accueillerait favorablement une affaire test\fg{} et utilisa la confusion qui suivit le verdict pour forcer le Congrès à légiférer \parencite{hazlett_rationality_1990}. Le \emph{New York Times} l'implora d'organiser un arrangement d'urgence avec l'industrie~; il s'y refusa.}. Le spectre ne fut pas régulé parce que le marché avait échoué~; il fut régulé parce qu'une solution de marché aurait privé les régulateurs et les incumbents de leurs rentes respectives.

\dimmark{D9}{*}%
Le second fait est la transformation du wireless en instrument politique de masse. La radiodiffusion, apparue aux États-Unis au début des années~1920 contre la volonté de l'industrie du point-à-point\footnote{\textcite{douglas_inventing_1987} a montré que le broadcasting naquit de la pratique des amateurs radio, pas d'une stratégie industrielle~: Marconi, les compagnies de câbles et le Post Office britannique étaient focalisés sur la communication point-à-point et ne voyaient pas la valeur de la diffusion à tous. La station KDKA de Pittsburgh (novembre~1920), propriété de Westinghouse, commença à émettre pour stimuler la vente de récepteurs, pas pour créer une industrie médiatique.}, créa le premier médium capable d'atteindre des millions de personnes simultanément. Les ventes de récepteurs radio aux États-Unis passèrent de 60~millions de dollars en~1922 à 850~millions en~1929 \parencite{kooij_information_2022}. RCA, sous la direction de David Sarnoff, vit son chiffre d'affaires bondir de 11~millions de dollars en~1922 à 60~millions en~1924 \parencite{kooij_information_2022}. La création de la National Broadcasting Company (NBC) en~1926, puis de la Columbia Broadcasting System (CBS) en~1928, organisa le broadcasting en réseaux nationaux financés par la publicité.

L'effet politique de ce nouveau médium a fait l'objet de deux études d'identification causale dont la convergence est remarquable. \textcite{adena_radio_2015}, exploitant la variation géographique et temporelle de la couverture radio en Allemagne entre~1929 et~1936 (signal calculé par le modèle de propagation ITM), ont montré que l'effet de la radio sur le soutien au parti nazi était bidirectionnel~: pendant la période démocratique (1929--1932), lorsque le contenu radiophonique était pro-Weimar, la radio ralentissait la progression du NSDAP dans les zones couvertes~; après la prise de contrôle du médium par les nazis en janvier~1933, la radio accélérait le soutien, le recrutement de membres et les actes antisémites. L'effet dépendait des prédispositions des auditeurs~: la propagande nazie était efficace dans les zones où l'antisémitisme historique était élevé, et contre-productive là où il était faible. Le même médium servait la démocratie ou la dictature selon qui le contrôlait~: le broadcasting est, dans la terminologie de la thèse, un pharmakon politique.

\dimmark{D9}{*}%
\textcite{gentzkow_television_2006}, par une stratégie d'identification similaire (variation géographique de l'introduction de la télévision aux États-Unis entre~1940 et~1960, exploitant la Seconde Guerre mondiale et le gel des licences FCC de~1948--1952 comme chocs exogènes), a montré que la télévision réduisait la participation électorale d'environ 2~points de pourcentage par décennie, expliquant la moitié du déclin de la participation aux élections non présidentielles depuis les années~1950. Le mécanisme identifié est le \emph{crowding out} de l'information politique~: la télévision, en tant que bien de divertissement supérieur, détournait les consommateurs des journaux et de la radio, qui contenaient davantage d'information politique locale. Le résultat est renforcé par un contraste frappant avec la radio~: \textcite{stromberg_radio_2004} a montré que la radio avait \emph{augmenté} la participation électorale dans les années~1930, probablement parce qu'elle causait moins de substitution à la lecture des journaux et apportait davantage d'information politique. Le broadcasting n'a donc pas un effet uniforme sur la participation~: la radio enrichit, la télévision appauvrit, et la différence tient à la structure de substitution entre les médias.

\emergence{D9~--- Le broadcasting crée le premier médium de masse. Son effet politique est bidirectionnel~: la radio renforce la démocratie (Weimar) ou la dictature (nazis) selon qui la contrôle \parencite{adena_radio_2015}~; la radio augmente la participation (Stromberg) là où la télévision la diminue (Gentzkow). Le même registre technique (TRANSMIT one-to-many) produit des effets opposés selon le contenu et la structure de substitution avec les autres médias. \questionch{2}{D9~--- Ce pharmakon politique est-il propre au broadcasting, ou se retrouve-t-il avec Internet et les réseaux sociaux~? Le mécanisme de crowding out de Gentzkow (la qualité du divertissement détourne de l'information) opère-t-il à nouveau avec le streaming et le smartphone~?}}


\subsection*{Bilan~: la communication sans fil au prisme des dix dimensions}\label{subsec:bilan_wireless}

Le tableau~\ref{tab:wireless-dimensions} récapitule ce que la traversée du cas wireless a permis d'établir. Trois niveaux de certitude sont distingués~: \textbf{*}~(résultat empirique solide, réplicable ou documenté indépendamment), \textbf{+}~(argument explicite mais fondation empirique incomplète), \textbf{?}~(piste ouverte).

\begin{table}[p]
\caption{La communication sans fil au prisme des dix dimensions~: observations, niveaux de certitude et questions ouvertes}\label{tab:wireless-dimensions}
\centering\footnotesize
\renewcommand{\arraystretch}{1.15}
\begin{tabularx}{\textwidth}{@{}cXcp{4.5cm}@{}}
\toprule
\textbf{Dim.} & \textbf{Observations du cas §\,\ref{sec:wireless}} & & \textbf{Question pour la suite} \\
\midrule
D2 & Trois mécanismes de monopole successifs~: brevet + leasing + refus d'interopérabilité (Marconi 1897--1914), licence étatique (Radio Act 1927), monopole créé par l'État (RCA 1919). Le spectre est le premier bien commun congestionnable immatériel de l'histoire de l'ICT. La régulation du spectre n'est pas une erreur mais un choix rationnel de distribution de rentes \parencite{hazlett_rationality_1990}. & * & La régulation du spectre préfigure-t-elle la régulation des plateformes numériques~? \\[3pt]
D5 & La Navy finance le réseau Pacifique (1,5~M\$), exproprie le wireless pendant la WWI, crée RCA. Même pattern que le Census pour Hollerith et la Convention pour le télégraphe Chappe~: l'État crée le marché avant que le marché n'existe. & * & L'État comme premier client est-il un fait constitutif de chaque vague d'ICT~? \\[3pt]
D8 & Données financières Marconi (ventes cumulées £6\,092 en~1900). RCA~: \$11M (1922) $\rightarrow$ \$60M (1924). Wireless bubble. Le pattern investissement $\rightarrow$ faillite $\rightarrow$ recapitalisation n'apparaît pas directement~; c'est l'expropriation étatique qui restructure le marché. & + & Le financement par la promesse stratégique (pas la rentabilité) est-il propre aux technologies de souveraineté~? \\[3pt]
D9 & Le broadcasting crée le premier médium de masse. Effet bidirectionnel~: radio pro-Weimar ralentit les nazis, radio nazie accélère le soutien \parencite{adena_radio_2015}. TV réduit la participation électorale \parencite{gentzkow_television_2006}, radio l'augmente \parencite{stromberg_radio_2004}. & * & Le pharmakon politique se reproduit-il avec Internet et les réseaux sociaux~? \\[3pt]
D10 & Le broadcasting naît contre la volonté de l'industrie du point-à-point. Les amateurs radio forcent le changement d'usage \parencite{douglas_inventing_1987}. Marconi, les câbles et le Post Office étaient focalisés sur le point-to-point. & + & L'user innovation comme mécanisme de changement de paradigme \\[3pt]
\midrule
0-TER & L'énergie passe du substrat (cuivre) au signal (onde). À efficacité spectrale fixée, la portée croît avec la racine carrée de la puissance d'émission~; doubler la portée exige de quadrupler la puissance. Mais cette loi est elle-même relâchée par l'efficacité de rayonnement~: 30~kW d'arc, concentrés sur une fréquence étroite, portent plus loin que 100~kW de spark dispersés en harmoniques \parencite{aitken_continuous_1985}. L'efficacité énergétique est la variable cachée de la trajectoire technique. C'est la première technologie ICT où le coût d'exploitation énergétique est continu et croissant avec l'usage. & * & Le déplacement de la matérialité du substrat au signal annonce-t-il le problème énergétique contemporain~? \\[3pt]
0-GOV & Les câbles sous-marins sont une vulnérabilité stratégique~: blocus câblé britannique, télégramme Zimmerman, interception commerciale \parencite{aitken_continuous_1985}. Le wireless naît comme technologie de souveraineté, pas de marché. La Navy orchestre la création de RCA. & * & La souveraineté numérique contemporaine prolonge-t-elle le même pattern~? \\[3pt]
\bottomrule
\end{tabularx}
\renewcommand{\arraystretch}{1.0}
\end{table}

La traversée du cas wireless suggère trois hypothèses que le chapitre~2 devra instruire. La première est que le spectre électromagnétique inaugure un type de ressource nouveau~: un bien commun congestionnable immatériel qui appelle une forme de régulation sans précédent dans l'histoire de l'ICT, et dont la régulation effective ne vise pas l'efficacité allocative mais la distribution de rentes. La seconde est que la transition du point-à-point au broadcasting est un changement de paradigme non anticipé par l'industrie en place, confirmant le pattern observé dans le cas du téléphone (§\,\ref{sec:telephone})~: les incumbents ne voient pas la valeur du nouveau registre parce que leur système métrique ne sait pas la mesurer. Marconi, comme Gray, ne manquait pas d'intelligence~; il manquait d'unités. La troisième hypothèse est propre à la dimension énergie-matière~: le wireless est la première technologie de l'information où le coût énergétique réside dans le signal lui-même, et non dans le substrat matériel du réseau. Le fil de cuivre coûtait cher à poser mais ne consommait rien une fois posé~; l'onde électromagnétique ne coûte rien à installer mais consomme en continu. Ce déplacement, du capital fixe matériel vers le flux énergétique d'exploitation, annonce le régime contemporain du datacenter~: une infrastructure dont le coût dominant n'est plus la construction mais l'alimentation.

\emergence{0-TER~--- Le wireless opère un basculement dans la structure de coûts de la transmission~: du capital fixe matériel (cuivre, gutta-percha) vers le flux énergétique d'exploitation (puissance d'émission). Ce basculement est masqué, dans les années~1900--1930, par les ordres de grandeur encore modestes~; mais la logique est en place, et elle se reproduira, amplifiée, à chaque étape de la numérisation. \questionch{3}{0-TER~--- Si le coût dominant de la transmission passe du substrat au signal, la trajectoire d'amélioration technique (spark $\rightarrow$ arc $\rightarrow$ tube $\rightarrow$ transistor) ne fait que repousser le moment où la consommation agrégée dépasse les gains d'efficacité unitaire. C'est le paradoxe de Jevons appliqué à la transmission.}}


